\documentclass[12pt, a4paper]{article}

%=============== PAQUETES ===============
\usepackage[utf8]{inputenc}
\usepackage[T1]{fontenc}
\usepackage[spanish,es-tabla]{babel}
\usepackage{geometry}
\usepackage{amsmath}
\usepackage{graphicx}
\usepackage{circuitikz}
\usetikzlibrary{babel}
\usepackage{siunitx}
\usepackage{hyperref}
\usepackage{enumitem}
\usepackage{booktabs}

%=============== CONFIGURACIÓN ===============
\geometry{a4paper, margin=2.5cm}

\hypersetup{
    colorlinks=true,
    linkcolor=blue,
    urlcolor=cyan,
    pdftitle={Guía Teórica Semana 1: Fundamentos de Circuitos DC},
}

% Configuración de unidades SI
\sisetup{
    output-decimal-marker = {.},
    per-mode = symbol
}

%=============== DOCUMENTO ===============
\begin{document}

\title{\textbf{Guía Teórica - Semana 1} \\ \large Fundamentos de Circuitos de Corriente Continua (DC)}
\author{Plan de Estudios de Electrónica}
\date{\today}
\maketitle

\begin{abstract}
La presente guía resume los conceptos fundamentales correspondientes a la Semana 1 del curso. Se abordan las magnitudes eléctricas básicas, las leyes fundamentales que rigen el comportamiento de los circuitos y las técnicas elementales para el análisis de circuitos puramente resistivos.
\end{abstract}

\tableofcontents
\newpage

\section{Conceptos Básicos}

\subsection{Magnitudes Eléctricas y Unidades}
Para el análisis de circuitos, es fundamental comprender las magnitudes físicas involucradas y sus unidades en el Sistema Internacional (SI):
\begin{itemize}
    \item \textbf{Carga Eléctrica ($q$):} Propiedad fundamental de la materia. Se mide en Culombios (C).
    \item \textbf{Corriente Eléctrica ($i$):} Es la tasa de flujo de carga eléctrica a través de una sección transversal. Se mide en Amperios (A). $i = \frac{dq}{dt}$
    \item \textbf{Tensión o Diferencia de Potencial ($v$):} Es la energía requerida para mover una carga unitaria a través de un elemento. Se mide en Voltios (V). $v = \frac{dw}{dq}$
    \item \textbf{Potencia ($p$):} Es la velocidad a la que se consume o se produce energía. Se mide en Vatios (W). $p = v \cdot i$
\end{itemize}

\subsection{Leyes de Kirchhoff}
Las leyes de Kirchhoff son los pilares del análisis de circuitos, basadas en la conservación de la carga y la energía.

\begin{description}
    \item[Ley de Corrientes de Kirchhoff (LKC):] La suma algebraica de las corrientes que entran a un nodo (punto de unión de dos o más elementos) es cero.
    \[ \sum_{n=1}^{N} i_n = 0 \]
    \item[Ley de Voltajes de Kirchhoff (LVK):] La suma algebraica de todas las caídas de tensión alrededor de una trayectoria cerrada (malla) es cero.
    \[ \sum_{m=1}^{M} v_m = 0 \]
\end{description}

\section{Circuitos Resistivos}

\subsection{La Resistencia y la Ley de Ohm}
La resistencia ($R$) es la oposición que presenta un material al flujo de corriente eléctrica, y se mide en Ohmios ($\Omega$). La relación fundamental entre tensión, corriente y resistencia está dada por la \textbf{Ley de Ohm}:
\[ v = i \cdot R \]

\subsection{Asociación de Resistencias y Reducción de Circuitos}

\subsubsection*{Resistencias en Serie}
Dos o más resistencias están en serie si comparten un único nodo y, por tanto, fluye la misma corriente a través de ellas. La resistencia equivalente es la suma de las resistencias individuales:
\[ R_{eq} = R_1 + R_2 + \dots + R_n \]

\subsubsection*{Resistencias en Paralelo}
Están en paralelo si están conectadas a los mismos dos nodos, por lo que tienen la misma tensión entre sus terminales. La inversa de la resistencia equivalente es la suma de las inversas:
\[ \frac{1}{R_{eq}} = \frac{1}{R_1} + \frac{1}{R_2} + \dots + \frac{1}{R_n} \]
Para dos resistencias en paralelo, la fórmula se simplifica a: $R_{eq} = \frac{R_1 \cdot R_2}{R_1 + R_2}$

\subsection{Divisores de Tensión y Corriente}

Son herramientas analíticas derivadas de las leyes de Kirchhoff y Ohm que permiten simplificar cálculos en circuitos básicos.

\begin{figure}[h!]
    \centering
    \begin{circuitikz}[scale=1.2]
        % Divisor de tensión
        \draw (0,0) to[V, l=$V_s$] (0,3)
              to[short] (2,3)
              to[R, l=$R_1$] (2,1.5)
              to[R, l=$R_2$, v=$V_2$] (2,0)
              to[short] (0,0);
        \node at (1,-0.5) {(a) Divisor de Tensión};
        
        % Divisor de corriente
        \draw (5,0) to[I, l=$I_s$] (5,3)
              to[short] (8,3);
        \draw (6.5,3) to[R, l=$R_1$, i=$i_1$] (6.5,0);
        \draw (8,3) to[R, l=$R_2$, i=$i_2$] (8,0);
        \draw (5,0) to[short] (8,0);
        \node at (6.5,-0.5) {(b) Divisor de Corriente};
    \end{circuitikz}
    \caption{Circuitos divisores fundamentales.}
\end{figure}

\begin{itemize}
    \item \textbf{Divisor de Tensión:} En un circuito serie, la tensión se divide proporcionalmente al valor de las resistencias. 
    \[ V_2 = V_s \cdot \frac{R_2}{R_1 + R_2} \]
    \item \textbf{Divisor de Corriente:} En un circuito paralelo, la corriente se divide de forma inversamente proporcional a las resistencias.
    \[ i_1 = I_s \cdot \frac{R_2}{R_1 + R_2} \quad \text{y} \quad i_2 = I_s \cdot \frac{R_1}{R_1 + R_2} \]
\end{itemize}

\section{Métodos de Análisis Formal}
Para circuitos que no pueden simplificarse únicamente con combinaciones serie/paralelo, se emplean métodos sistemáticos.

\subsection{Análisis Nodal}
Basado en la Ley de Corrientes de Kirchhoff (LKC). El procedimiento básico es:
\begin{enumerate}[noitemsep]
    \item Seleccionar un nodo de referencia (tierra, $0$ V).
    \item Asignar variables de tensión a los nodos restantes ($v_1, v_2, \dots$).
    \item Aplicar la LKC en cada nodo desconocido, expresando las corrientes en función de las tensiones nodales usando la Ley de Ohm.
    \item Resolver el sistema de ecuaciones algebraicas resultante.
\end{enumerate}

\subsection{Análisis de Mallas}
Basado en la Ley de Voltajes de Kirchhoff (LVK), aplicable a circuitos planos (que pueden dibujarse sin cruces de ramas). El procedimiento es:
\begin{enumerate}[noitemsep]
    \item Identificar las mallas del circuito.
    \item Asignar corrientes de malla ($i_1, i_2, \dots$), preferiblemente en sentido horario.
    \item Aplicar la LVK a cada malla, expresando las tensiones en función de las corrientes de malla mediante la Ley de Ohm.
    \item Resolver el sistema de ecuaciones para encontrar las corrientes de malla.
\end{enumerate}

\newpage
\section{Ejercicios Resueltos}

\subsection*{Ejercicio 1: Ley de Ohm}
\textbf{Problema:} Un resistor de \SI{4}{\ohm} está conectado a una fuente de \SI{12}{\volt}. Calcule la corriente que fluye por el resistor y la potencia disipada. \\
\textbf{Solución:} 
Aplicando la Ley de Ohm: $i = \frac{v}{R} = \frac{12}{4} = \SI{3}{\ampere}$.\\
La potencia disipada es: $p = v \cdot i = 12 \cdot 3 = \SI{36}{\watt}$.

\subsection*{Ejercicio 2: Resistencias en Serie}
\textbf{Problema:} Tres resistores de \SI{10}{\ohm}, \SI{20}{\ohm} y \SI{30}{\ohm} están conectados en serie. ¿Cuál es la resistencia equivalente? \\
\textbf{Solución:} 
Para resistores en serie, se suman sus valores:
$R_{eq} = R_1 + R_2 + R_3 = 10 + 20 + 30 = \SI{60}{\ohm}$.

\subsection*{Ejercicio 3: Resistencias en Paralelo}
\textbf{Problema:} Dos resistores de \SI{100}{\ohm} y \SI{25}{\ohm} se conectan en paralelo. Determine la resistencia equivalente.\\
\textbf{Solución:} 
Usando la fórmula para dos resistores en paralelo:
$R_{eq} = \frac{R_1 \cdot R_2}{R_1 + R_2} = \frac{100 \cdot 25}{100 + 25} = \frac{2500}{125} = \SI{20}{\ohm}$.

\subsection*{Ejercicio 4: Reducción Mixta de Circuito}
\textbf{Problema:} Un resistor $R_1 = \SI{10}{\ohm}$ está en serie con una combinación en paralelo de $R_2 = \SI{30}{\ohm}$ y $R_3 = \SI{60}{\ohm}$. Calcule la resistencia total.\\
\begin{figure}[h!]
    \centering
    \begin{circuitikz}[scale=0.9]
        \draw (0,0) to[V, l=$V_s$] (0,4)
              to[R, l={$R_1$}] (3,4)
              to[short] (3,3)
              to[R, l={$R_2$}] (3,1)
              to[short] (3,0) to[short] (0,0);
        \draw (3,3) to[short] (5,3) to[R, l={$R_3$}] (5,1) to[short] (3,1);
    \end{circuitikz}
    \caption{Circuito mixto para el Ejercicio 4.}
\end{figure}
\textbf{Solución:} 
Primero resolvemos el paralelo: $R_{23} = \frac{30 \cdot 60}{30 + 60} = \frac{1800}{90} = \SI{20}{\ohm}$.\\
Luego, sumamos la resistencia en serie: $R_{eq} = R_1 + R_{23} = 10 + 20 = \SI{30}{\ohm}$.

\subsection*{Ejercicio 5: Divisor de Tensión}
\textbf{Problema:} Dos resistores $R_1 = \SI{2}{\kilo\ohm}$ y $R_2 = \SI{6}{\kilo\ohm}$ están en serie conectados a una fuente de \SI{24}{\volt}. ¿Cuál es la tensión en los terminales de $R_2$? \\
\begin{figure}[h!]
    \centering
    \begin{circuitikz}
        \draw (0,0) to[V, l={\SI{24}{\volt}}] (0,3)
              to[R, l={$R_1$}] (3,3)
              to[R, l={$R_2$}, v=$V_2$] (6,3)
              to[short] (6,0) to[short] (0,0);
    \end{circuitikz}
    \caption{Divisor de tensión para el Ejercicio 5.}
\end{figure}
\textbf{Solución:} 
Usando la fórmula del divisor de tensión:
$V_2 = V_s \frac{R_2}{R_1 + R_2} = 24 \frac{6}{2 + 6} = 24 \left(\frac{6}{8}\right) = \SI{18}{\volt}$.

\subsection*{Ejercicio 6: Divisor de Corriente}
\textbf{Problema:} Una fuente de corriente ideal de \SI{10}{\milli\ampere} alimenta un circuito con dos ramas en paralelo, $R_1 = \SI{1}{\kilo\ohm}$ y $R_2 = \SI{4}{\kilo\ohm}$. Determine la corriente que atraviesa $R_1$.\\
\begin{figure}[h!]
    \centering
    \begin{circuitikz}
        \draw (0,0) to[I, l={\SI{10}{\milli\ampere}}] (0,3)
              to[short] (4,3);
        \draw (1.5,3) to[R, l={$R_1$}, i=$i_1$] (1.5,0);
        \draw (4,3) to[R, l={$R_2$}] (4,0);
        \draw (0,0) to[short] (4,0);
    \end{circuitikz}
    \caption{Divisor de corriente para el Ejercicio 6.}
\end{figure}
\textbf{Solución:} 
Aplicando el divisor de corriente (note que el numerador es la \emph{otra} resistencia):
$i_1 = I_s \frac{R_2}{R_1 + R_2} = 10 \frac{4}{1 + 4} = 10 \left(\frac{4}{5}\right) = \SI{8}{\milli\ampere}$.

\subsection*{Ejercicio 7: Ley de Corrientes de Kirchhoff (LKC)}
\textbf{Problema:} En un nodo de un circuito concurren tres ramas. La corriente $i_1 = \SI{5}{\ampere}$ entra al nodo y la corriente $i_2 = \SI{2}{\ampere}$ sale del nodo. ¿Cuál es el valor y dirección de la tercera corriente $i_3$? \\
\textbf{Solución:} 
Definiendo corrientes entrantes como positivas y salientes como negativas:
$\sum i = 0 \implies i_1 - i_2 - i_3 = 0 \implies 5 - 2 - i_3 = 0 \implies i_3 = \SI{3}{\ampere}$.\\
Como $i_3$ es positiva asumiendo que "sale", la corriente $i_3$ sale del nodo con un valor de \SI{3}{\ampere}.

\subsection*{Ejercicio 8: Ley de Voltajes de Kirchhoff (LVK)}
\textbf{Problema:} En un lazo cerrado hay una fuente de tensión $V_s = \SI{15}{\volt}$ y tres resistores en serie. Las caídas de tensión en los primeros dos resistores son $V_1 = \SI{4}{\volt}$ y $V_2 = \SI{6}{\volt}$. Calcule la caída de tensión en el tercer resistor ($V_3$).\\
\textbf{Solución:} 
Recorriendo la malla y aplicando la LVK:
$V_s - V_1 - V_2 - V_3 = 0 \implies 15 - 4 - 6 - V_3 = 0 \implies V_3 = \SI{5}{\volt}$.

\subsection*{Ejercicio 9: Potencia y Energía}
\textbf{Problema:} Un bombillo tiene una resistencia de \SI{240}{\ohm} y se conecta a una red de \SI{120}{\volt}. Calcule la potencia consumida y la energía disipada (en julios) si se mantiene encendido durante 1 hora.\\
\textbf{Solución:} 
Potencia: $p = \frac{v^2}{R} = \frac{120^2}{240} = \frac{14400}{240} = \SI{60}{\watt}$.\\
Energía: $w = p \cdot t$. El tiempo en segundos es $t = 1 \text{ h} \cdot 3600 \text{ s/h} = \SI{3600}{\second}$.\\
$w = 60 \cdot 3600 = \SI{216000}{\joule} = \SI{216}{\kilo\joule}$.

\subsection*{Ejercicio 10: Análisis de Nodos Básico}
\textbf{Problema:} Un nodo $A$ tiene un voltaje $v_A$ desconocido. Está conectado a un nodo de \SI{10}{\volt} a través de una resistencia de \SI{2}{\ohm}, y a tierra (\SI{0}{\volt}) a través de una resistencia de \SI{3}{\ohm}. Una fuente de corriente inyecta \SI{2}{\ampere} directamente en el nodo $A$. Escriba la ecuación del nodo y halle $v_A$.\\
\begin{figure}[h!]
    \centering
    \begin{circuitikz}
        % Línea de tierra
        \draw (0,0) -- (4,0);
        \node[ground] at (2,0) {};
        
        % Rama izquierda: fuente de 10V y resistencia de 2 ohmios
        \draw (0,0) to[V, l={\SI{10}{\volt}}] (0,2)
              to[R] (2,2);
              
        % Rama central: resistencia de 3 ohmios a tierra
        \draw (2,2) to[R] (2,0);
        
        % Rama derecha: fuente de corriente de 2A (inyectando al nodo A)
        \draw (4,0) to[I, l={\SI{2}{\ampere}}] (4,2)
              to[short] (2,2);
              
        % Etiqueta del Nodo A
        \node[circ] at (2,2) {};
        \node[above] at (2,2.2) {Nodo $A$ ($v_A$)};
    \end{circuitikz}
    \caption{Circuito nodal para el Ejercicio 10.}
\end{figure}
\textbf{Solución:} 
Aplicamos la LKC en el nodo $A$, asumiendo que todas las corrientes resistivas salen del nodo:
\[ \frac{v_A - 10}{2} + \frac{v_A - 0}{3} - 2 = 0 \]
Multiplicando todo por 6 (mínimo común múltiplo):
\[ 3(v_A - 10) + 2(v_A) - 12 = 0 \]
\[ 3v_A - 30 + 2v_A - 12 = 0 \]
\[ 5v_A = 42 \implies v_A = \SI{8.4}{\volt} \]

\vspace{1cm}
\noindent\rule{\textwidth}{0.4pt}
\textbf{Recomendación de Estudio:} Repasar los ejemplos resueltos de los capítulos 1 y 2 del texto base y practicar dibujando e identificando nodos, ramas y mallas en los simuladores de circuitos (ej. DroidTesla, LTspice).

\end{document}